% Shared house style for the workflow schematics.
%
% Both supplementary workflow figures are the same kind of diagram: a titled
% header, a numbered chain of steps that wraps onto a second row, and a closing
% note. Everything about how that looks lives here, so the two figure files are
% content only and cannot drift apart visually.
%
% Included by the figure files; not a standalone document.

\usepackage[T1]{fontenc}
\usepackage{lmodern}
\usepackage{tikz}
\usetikzlibrary{arrows.meta, positioning, calc}

% --- palette ---------------------------------------------------------------
% One accent per step role, dark enough to stay legible in greyscale. Header
% bars carry the accent at full strength; bodies carry an 8% tint of it.
\definecolor{wfInk}{HTML}{1F2D3D}      % body text
\definecolor{wfMuted}{HTML}{5A6B7B}    % secondary text, arrows
\definecolor{wfRule}{HTML}{C7D0D9}
\definecolor{wfSetup}{HTML}{B07D2B}    % scope-setting steps
\definecolor{wfFit}{HTML}{2E6DA4}      % fitting and prediction steps
\definecolor{wfInner}{HTML}{2E7D52}    % nested-resampling steps

% --- geometry --------------------------------------------------------------
% One box width drives everything; the header and body derive their text widths
% from it so the two nodes are exactly the same total width and align on their
% north-west corners.
\newlength{\wfBoxW}   \setlength{\wfBoxW}{52mm}
\newlength{\wfBoxH}   \setlength{\wfBoxH}{42mm}
\newlength{\wfPad}    \setlength{\wfPad}{2.6mm}
\newlength{\wfTextW}  \setlength{\wfTextW}{\dimexpr\wfBoxW-2\wfPad\relax}
\newlength{\wfHeadH}  \setlength{\wfHeadH}{7.4mm}
\newlength{\wfGapX}   \setlength{\wfGapX}{7mm}
% Five boxes and four gaps: the row-1 chain sets the canvas width, which the
% title and footnote blocks then span.
\newlength{\wfCanvasW}\setlength{\wfCanvasW}{\dimexpr 5\wfBoxW+4\wfGapX\relax}

% Resize the step boxes. The tikz styles read \wfTextW and \wfBoxH when a node
% is created, so setting these before a \wfstep call changes that box and every
% later one until it is set again. Row 2 uses wider, shorter boxes because its
% steps carry more text and there are only two of them.
\newcommand{\wfsetbox}[2]{%
  \setlength{\wfBoxW}{#1}%
  \setlength{\wfBoxH}{#2}%
  \setlength{\wfTextW}{\dimexpr\wfBoxW-2\wfPad\relax}%
}

\tikzset{
  wfbody/.style={
    draw=#1, fill=#1!8, line width=0.5pt, rounded corners=1.8pt,
    text width=\wfTextW, minimum height=\wfBoxH, inner xsep=\wfPad,
    inner ysep=0pt, align=left, anchor=north west,
  },
  wfhead/.style={
    fill=#1, text=white, text width=\wfTextW, inner xsep=\wfPad,
    inner ysep=1.9mm, align=left, anchor=north west, rounded corners=1.8pt,
    font=\sffamily\bfseries\small,
  },
  wfarrow/.style={
    -{Stealth[length=2.4mm, width=1.7mm]}, draw=wfMuted, line width=0.8pt,
  },
}

% \wfstep{<accent>}{<node name>}{<placement>}{<number>}{<title>}{<body>}
%
% The body node is drawn first and reserves the header's height at the top; the
% header bar is then laid over that reserved strip. Body text is a text width,
% so paragraphs wrap themselves and every box keeps the same footprint no matter
% how much text a step carries.
% The header height is reserved with a zero-width strut rather than \vspace*:
% inside a TikZ node with align=left the content starts in horizontal mode, so
% a leading \vspace* is applied after the first line has already been set and
% the header then covers that line. A strut forces a line box of exactly the
% header's height, whatever mode we are in.
\newcommand{\wfstep}[6]{%
  \node[wfbody=#1, #3] (#2) {%
    \rule{0pt}{\dimexpr\wfHeadH+2.2mm\relax}\par
    \color{wfInk}\sffamily\fontsize{7.6}{10}\selectfont\raggedright #6%
  };%
  \node[wfhead=#1] at (#2.north west) {#4\hspace{1.4ex}#5};%
}

% A step body is a list of short statements; \wfsep separates them.
\newcommand{\wfsep}{\par\vspace{1.5mm}}

% \wftitle{<title>}{<subtitle>}
\newcommand{\wftitle}[2]{%
  {\color{wfInk}\sffamily\bfseries\fontsize{17}{20}\selectfont #1\par}%
  \vspace{1.8mm}%
  {\color{wfMuted}\sffamily\fontsize{9}{11}\selectfont #2\par}%
  \vspace{2.8mm}%
  {\color{wfRule}\rule{\linewidth}{0.6pt}}%
}

% \wffoot{<note>}
\newcommand{\wffoot}[1]{%
  {\color{wfRule}\rule{\linewidth}{0.6pt}}\par\vspace{1.9mm}%
  {\color{wfMuted}\sffamily\itshape\fontsize{7.8}{9.6}\selectfont #1\par}%
}
